\chapter{Introduction}

The aim of this work is twofold: to give a foundational account of comonadic techniques in unstable homotopy theory, and to apply this foundation to prove an unstable variant of the Morava change of rings theorem from stable homotopy theory \cite{mr,greenbook}, completely formalized in Lean~4 \cite{lean4,mathlib}.

Comonads have traditionally been used to define resolving objects, and hence derived functors, in both abelian and non-abelian settings \cite{barrbeck}. In work of Bousfield \cite{bousfield}, a comonad $G$ defines a model structure on the category of cosimplicial objects over the category $\mathcal{C}^G$ of coalgebras over $G$, in which the traditional cobar resolutions compute derived mapping spaces and $\Ext$ groups.

Under suitable hypotheses, a morphism of comonads induces a functor on the categories of cosimplicial coalgebras that preserves resolving objects. This is the content of our \emph{Generalized Shapiro's Lemma for Comonads} (Proposition~\ref{prop:generalized_shapiro}). We then define an unstable version $U_{\Sigma}$ of the Morava stabilizer algebra $\Sigma(n)$ and use the generalized Shapiro's lemma to prove that the unstable change of rings isomorphism holds in connectivity $k \ge 2p^n + 1$ (Theorem~\ref{thm:unstable_change_of_rings}).

\chapter{The Theory of Comonads}

\section{Comonads, Coalgebras, and the Cokleisli Category}

\begin{definition}[Cokleisli-to-Coalgebra Functor]
  \label{defn:cokleisli_to_coalgebra}
  \lean{ThesisLean.Section2_1.cokleisliToCoalgebra}
  \leanok
  For a comonad $G$ on a category $\mathcal{C}$, the functor $i : \mathcal{C}_G \to \mathcal{C}^G$ sends $X \mapsto (GX, \Delta_X)$ and a Kleisli morphism $f : GX \to Y$ to $G(f) \circ \Delta_X : GX \to GY$.
\end{definition}

\begin{definition}[Coalgebra-to-Cokleisli Functor]
  \label{defn:coalgebra_to_cokleisli}
  \lean{ThesisLean.Section2_1.coalgebraToCokleisli}
  \leanok
  For a comonad $G$ on $\mathcal{C}$, the functor $j : \mathcal{C}^G \to \mathcal{C}_G$ sends a $G$-coalgebra $(X, \psi_X)$ to $X$ and a coalgebra morphism $f : X \to Y$ to $f \circ \epsilon_X : GX \to Y$.
\end{definition}

\begin{proposition}[Cokleisli--Coalgebra Comparison]
  \label{prop:cokleisli_adj}
  \lean{ThesisLean.Section2_1.cokleisliToCoalgebraFullyFaithful}
  \uses{defn:cokleisli_to_coalgebra, defn:coalgebra_to_cokleisli}
  \leanok
  The functor $i : \mathcal{C}_G \to \mathcal{C}^G$ is fully faithful (exhibiting $\mathcal{C}_G$ as the full subcategory of cofree coalgebras in $\mathcal{C}^G$), and $j : \mathcal{C}^G \to \mathcal{C}_G$ provides the comparison functor in the opposite direction.
\end{proposition}
\begin{proof}
  \leanok
  Verified via the hom-set equivalence $\mathcal{C}(GX, Y) \simeq \mathcal{C}^G((X,\psi_X), (GY,\Delta_Y))$ sending $f \mapsto G(f) \circ \Delta_X \circ \psi_X$ with inverse $g \mapsto \epsilon_Y \circ g$.
\end{proof}

\section{Morphisms of Comonads and Induced Functors}

\begin{definition}[Functor of Comonads]
  \label{defn:comonad_functor}
  \lean{ThesisLean.Section2_2.ComonadFunctor}
  \leanok
  A functor of comonads $(F, \alpha) : (\mathcal{C}, G) \to (\mathcal{D}, H)$ consists of a functor $F : \mathcal{C} \to \mathcal{D}$ and a natural transformation $\alpha : F \circ G \to H \circ F$ satisfying $\Delta^H_{FX} \circ \alpha_X = H(\alpha_X) \circ \alpha_{GX} \circ F(\Delta^G_X)$ and $\epsilon^H_{FX} \circ \alpha_X = F(\epsilon^G_X)$.
\end{definition}

\begin{proposition}[Coalgebra Pushforward]
  \label{prop:comonad_pushforward}
  \lean{ThesisLean.Section2_2.ComonadFunctor.pushforward}
  \uses{defn:comonad_functor}
  \leanok
  A functor of comonads $(F, \alpha) : (\mathcal{C}, G) \to (\mathcal{D}, H)$ induces a pushforward functor $\alpha_* : \mathcal{C}^G \to \mathcal{D}^H$ given on objects by $\alpha_*(X, \psi_X) = (FX, \alpha_X \circ F(\psi_X))$, commuting strictly with the forgetful functors.
\end{proposition}
\begin{proof}
  \leanok
  Coassociativity and counitality of $\alpha_X \circ F(\psi_X)$ follow directly from the coalgebra axioms for $(X, \psi_X)$, naturality of $\alpha$, and Definition~\ref{defn:comonad_functor}.
\end{proof}

\begin{definition}[Opfunctor of Comonads]
  \label{defn:comonad_opfunctor}
  \lean{ThesisLean.Section2_2.ComonadOpfunctor}
  \leanok
  An opfunctor of comonads $(U, \beta) : (\mathcal{D}, H) \to (\mathcal{C}, G)$ consists of a functor $U : \mathcal{D} \to \mathcal{C}$ and a natural transformation $\beta : G \circ U \to U \circ H$ satisfying $U(\Delta^H_Y) \circ \beta_Y = \beta_{HY} \circ G(\beta_Y) \circ \Delta^G_{UY}$ and $U(\epsilon^H_Y) \circ \beta_Y = \epsilon^G_{UY}$.
\end{definition}

\begin{proposition}[Cokleisli Pushforward]
  \label{prop:cokleisli_pushforward}
  \lean{ThesisLean.Section2_2.ComonadOpfunctor.cokleisliPushforward}
  \uses{defn:comonad_opfunctor}
  \leanok
  An opfunctor of comonads $(U, \beta) : (\mathcal{D}, H) \to (\mathcal{C}, G)$ induces a functor $\beta_* : \mathcal{D}_H \to \mathcal{C}_G$ given by $\beta_*(Y) = UY$ and $\beta_*(f : HY \to Y') = U(f) \circ \beta_Y$.
\end{proposition}
\begin{proof}
  \leanok
  Preservation of identities and composition follows from the two opfunctor axioms of Definition~\ref{defn:comonad_opfunctor}.
\end{proof}

\begin{proposition}[Coalgebra Pullback via Equalizers]
  \label{prop:coalgebra_pullback}
  \lean{ThesisLean.Section2_2.ComonadOpfunctor.coalgebraPullback}
  \uses{defn:comonad_opfunctor}
  \leanok
  Assume $\mathcal{C}^G$ has equalizers. Any opfunctor of comonads $(U, \beta) : (\mathcal{D}, H) \to (\mathcal{C}, G)$ induces a pullback functor $\alpha^* : \mathcal{D}^H \to \mathcal{C}^G$ defined on $(Y, \psi_Y)$ as the equalizer in $\mathcal{C}^G$ of $\cofree_G(U\psi_Y)$ and $G(\beta_Y) \circ \Delta^G_{UY} : \cofree_G(UY) \rightrightarrows \cofree_G(UHY)$.
\end{proposition}
\begin{proof}
  \leanok
  Constructed via the limit equalizer of the parallel pair of cofree $G$-coalgebra morphisms.
\end{proof}

\section{The Mate Correspondence and Lifted Adjunctions}

\begin{proposition}[Fundamental Mate Relations]
  \label{prop:mate_relations}
  \lean{ThesisLean.Section2_3.mate_fundamental_rel_counit, ThesisLean.Section2_3.mate_fundamental_rel_unit}
  \uses{defn:comonad_functor}
  \leanok
  Given $F \dashv U$ with unit $\sigma : 1_{\mathcal{C}} \to UF$ and counit $\tau : FU \to 1_{\mathcal{D}}$, and a comonad functor $(F, \alpha) : (\mathcal{C}, G) \to (\mathcal{D}, H)$, the mate transformation $\beta_Y := U(H(\tau_Y) \circ \alpha_{UY}) \circ \sigma_{GUY}$ satisfies
  \[
    \tau_{HY} \circ F(\beta_Y) = H(\tau_Y) \circ \alpha_{UY} \quad\text{and}\quad \beta_{FX} \circ G(\sigma_X) = U(\alpha_X) \circ \sigma_{GX}.
  \]
\end{proposition}
\begin{proof}
  \leanok
  Proved using naturality of $\tau, \sigma, \alpha$ and the triangle identities for $F \dashv U$.
\end{proof}

\begin{proposition}[Mate Opfunctor of Comonads]
  \label{prop:mate_opfunctor}
  \lean{ThesisLean.Section2_3.mateOpfunctor}
  \uses{defn:comonad_functor, defn:comonad_opfunctor, prop:mate_relations}
  \leanok
  Under the hypotheses of Proposition~\ref{prop:mate_relations}, $(U, \beta)$ is an opfunctor of comonads $(\mathcal{D}, H) \to (\mathcal{C}, G)$.
\end{proposition}
\begin{proof}
  \leanok
  Verified by passing to left-adjoint transposes via Proposition~\ref{prop:mate_relations} and applying the comonad functor laws of $(F, \alpha)$.
\end{proof}

\begin{proposition}[Lifted Adjunction on Coalgebras]
  \label{prop:lifted_adjunction}
  \lean{ThesisLean.Section2_3.liftedAdjunction}
  \uses{prop:comonad_pushforward, prop:coalgebra_pullback, prop:mate_opfunctor}
  \leanok
  When $\mathcal{C}^G$ has equalizers, the pushforward $\alpha_* : \mathcal{C}^G \to \mathcal{D}^H$ is left adjoint to the pullback $\alpha^* : \mathcal{D}^H \to \mathcal{C}^G$.
\end{proposition}
\begin{proof}
  \leanok
  Constructed by lifting the unit $\sigma_X$ into the equalizer $\alpha^*\alpha_* X$ via $G(\sigma_X) \circ \psi_X$, lifting the counit $\tau_Y$ via $\tau_Y \circ F(\epsilon^G_{UY} \circ \iota_Y)$, and verifying both triangle identities in $\mathcal{D}^H$ and $\mathcal{C}^G$.
\end{proof}

\section{The Cobar Resolution and Left Contractibility}

\begin{definition}[Induced Monad and Cosimplicial Cobar Complex]
  \label{defn:cobar_complex}
  \lean{ThesisLean.Section2_4.inducedMonadOnCoalgebras, ThesisLean.Section2_4.cobarObj, ThesisLean.Section2_4.cobarCoface, ThesisLean.Section2_4.cobarCodegeneracy}
  \leanok
  For a comonad $G$ on $\mathcal{C}$, the adjunction $U_G \dashv \cofree_G$ induces a monad $(T, \mu, \eta)$ on $\mathcal{C}^G$ with $T(X, \psi_X) = (GX, \Delta_X)$, $\eta_X = \psi_X$, and $\mu_X = G(\epsilon_X)$. The cosimplicial cobar complex $C^{\bullet}_G(X)$ has $C^n_G(X) = T^{n+1}(X)$ with coaugmentation $d^0_{-1} = \eta_X$, cofaces $d^i_n : T^{n+1}X \to T^{n+2}X$, and codegeneracies $s^j_n : T^{n+2}X \to T^{n+1}X$.
\end{definition}

\begin{proposition}[Cosimplicial Identities and Left Contractibility into Cofree Coalgebras]
  \label{prop:cobar_contractible}
  \lean{ThesisLean.Section2_4.cobarCoface_comp_cobarCoface, ThesisLean.Section2_4.cobarHomLeftContractibleTower}
  \uses{defn:cobar_complex}
  \leanok
  The cobar complex $C^{\bullet}_G(X)$ satisfies all five cosimplicial identities, and for every $J \in \mathcal{C}$ the augmented simplicial hom-tower $\Hom_{\mathcal{C}^G}(C^{\bullet}_G(X), G(J)) \to \Hom_{\mathcal{C}^G}(X, G(J))$ admits a left contraction given by $f \mapsto G(f) \circ \Delta$.
\end{proposition}
\begin{proof}
  \leanok
  All five cosimplicial identities follow by induction on cosimplicial degree from the monad laws of $T$. The left contraction identities follow from coassociativity and counitality of $G$.
\end{proof}

\section{A Generalized Shapiro's Lemma for Comonads}

\begin{proposition}[Shapiro Naturality Triangle]
  \label{prop:shapiro_triangle}
  \lean{ThesisLean.Section2_5.shapiro_naturality_triangle_coaug}
  \uses{prop:lifted_adjunction, defn:cobar_complex}
  \leanok
  For any $f : M \to \alpha^* N$ in $\mathcal{C}^G$, the lifted adjunction $\tilde{\gamma} : \Hom_{\mathcal{C}^G}(M, \alpha^* -) \simeq \Hom_{\mathcal{D}^H}(\alpha_* M, -)$ commutes with the cobar coaugmentation $\psi_N : N \to H(N)$:
  \[
    \tilde{\gamma}_{M, H(N)}(f \circ_{\text{post}} \alpha^*(\psi_N)) = \tilde{\gamma}_{M, N}(f) \circ \psi_N.
  \]
\end{proposition}
\begin{proof}
  \leanok
  Proved via naturality of the lifted unit and the adjunction hom-equivalence.
\end{proof}

\begin{proposition}[Generalized Shapiro's Lemma for Comonads]
  \label{prop:generalized_shapiro}
  \lean{ThesisLean.Section2_5.pullbackCobarCofreeLeftContractibleTower, ThesisLean.Section2_5.generalizedShapiroCobarIso, ThesisLean.Section2_5.generalizedShapiroExtCobarIso}
  \uses{prop:lifted_adjunction, prop:cobar_contractible, prop:shapiro_triangle}
  \leanok
  Assume there is a functor $W : \mathcal{D} \to \mathcal{C}$ and a natural isomorphism $U_G \circ \alpha^* \cong W \circ U_H$. Then:
  \begin{enumerate}
    \item For every $N \in \mathcal{D}^H$ and $J \in \mathcal{C}$, the hom-tower $\Hom_{\mathcal{C}^G}(\alpha^*(C^{\bullet}_H(N)), G(J)) \to \Hom_{\mathcal{C}^G}(\alpha^* N, G(J))$ is left-contractible, so $\alpha^* N \to \alpha^*(C^{\bullet}_H(N))$ is a Bousfield weak $G$-resolution.
    \item For every $M \in \mathcal{C}^G$ and $N \in \mathcal{D}^H$, there is an isomorphism of coaugmented cosimplicial hom-towers $\Hom_{\mathcal{C}^G}(M, \alpha^*(C^{\bullet}_H(N))) \cong \Hom_{\mathcal{D}^H}(\alpha_* M, C^{\bullet}_H(N))$, inducing $\Ext_G(M, \alpha^* N) \cong \Ext_H(\alpha_* M, N)$.
  \end{enumerate}
\end{proposition}
\begin{proof}
  \leanok
  Part (1) transports the underlying left contraction $W(\epsilon^H)^*$ across the adjunction $U_G \dashv \cofree_G$ and the isomorphism $U_G \circ \alpha^* \cong W \circ U_H$. Part (2) applies the lifted adjunction levelwise across the cobar complex.
\end{proof}

\chapter{An Unstable Change of Rings}

\section{The Stable Case and Canonical Coalgebra Equalizer}

\begin{proposition}[Canonical Equalizer Presentation of Coalgebras]
  \label{prop:canonical_equalizer}
  \lean{ThesisLean.Section3_1.coalgebraCanonicalEqualizerIso}
  \leanok
  For any comonad $G$ on $\mathcal{C}$ and any $G$-coalgebra $X = (X, \psi_X)$, the coaction $\psi_X : X \to \cofree_G(X)$ exhibits $X$ as the equalizer in $\mathcal{C}^G$ of $\cofree_G(\psi_X)$ and $\Delta_X : \cofree_G(X) \rightrightarrows \cofree_G(GX)$.
\end{proposition}
\begin{proof}
  \leanok
  Any cone $k : W \to \cofree_G(X)$ equalizing $\cofree_G(\psi_X)$ and $\Delta_X$ factors uniquely through $X$ via $\epsilon_X \circ k$.
\end{proof}

\begin{theorem}[Stable Change of Rings]
  \label{thm:stable_cor}
  \lean{ThesisLean.Section3_1.stableChangeOfRingsCobarIso}
  \uses{prop:generalized_shapiro}
  \leanok
  When the stable cotensor satisfies $U_G \circ \pi^* \cong (B(n)_* \otimes_{K(n)_*} -) \circ U_H$, Proposition~\ref{prop:generalized_shapiro} yields the isomorphism of cobar hom-towers $\Hom_{\mathcal{C}^G}(M, \pi^* C^{\bullet}_H(N)) \cong \Hom_{\mathcal{D}^H}(\pi_* M, C^{\bullet}_H(N))$.
\end{theorem}
\begin{proof}
  \leanok
  Direct specialization of Proposition~\ref{prop:generalized_shapiro}.
\end{proof}

\section{The Unstable Case}

\begin{theorem}[Unstable Sandwich Comonad $U_{\Sigma}$]
  \label{thm:unstable_sigma_comonad}
  \lean{ThesisLean.Section3_2.unstableSigmaComonad}
  \leanok
  For any comonad $U_{\Gamma}$ on $\mathcal{C}$ and any adjunction $F \dashv U : \mathcal{D} \to \mathcal{C}$ with unit $\eta$ and counit $\varepsilon$, the endofunctor $U_{\Sigma} := F \circ U_{\Gamma} \circ U : \mathcal{D} \to \mathcal{D}$ is a comonad with counit $\epsilon^{\Sigma}_N = \varepsilon_N \circ F(\epsilon^{\Gamma}_{UN})$ and comultiplication $\Delta^{\Sigma}_N = F(U_{\Gamma}(\eta_{U_{\Gamma} UN})) \circ F(\Delta^{\Gamma}_{UN})$.
\end{theorem}
\begin{proof}
  \leanok
  Coassociativity follows from naturality of $\Delta^{\Gamma}$, coassociativity of $U_{\Gamma}$, and naturality of $\eta$. Counitality follows from the counit laws of $U_{\Gamma}$ and both triangle identities of $F \dashv U$.
\end{proof}

\begin{theorem}[Unstable Base-Change Comonad Morphism]
  \label{thm:unstable_alpha}
  \lean{ThesisLean.Section3_2.unstableComonadFunctor}
  \uses{thm:unstable_sigma_comonad, defn:comonad_functor}
  \leanok
  The natural transformation $\alpha_M := F(U_{\Gamma}(\eta_M)) : F(U_{\Gamma}(M)) \to U_{\Sigma}(F(M))$ makes $(F, \alpha) : (\mathcal{C}, U_{\Gamma}) \to (\mathcal{D}, U_{\Sigma})$ a functor of comonads.
\end{theorem}
\begin{proof}
  \leanok
  Verified via naturality of $\Delta^{\Gamma}$, $\epsilon^{\Gamma}$, and $\eta$, together with the left triangle identity of $F \dashv U$.
\end{proof}

\begin{theorem}[Unstable Lifted Adjunction $\alpha_* \dashv \alpha^*$]
  \label{thm:unstable_lifted_adj}
  \lean{ThesisLean.Section3_2.unstableLiftedAdjunction, ThesisLean.Section3_2.unstableMateNatTrans_app}
  \uses{thm:unstable_alpha, prop:lifted_adjunction}
  \leanok
  When $\mathcal{C}^{U_{\Gamma}}$ has equalizers, $(F, \alpha)$ induces an adjoint pair $\alpha_* \dashv \alpha^* : \mathcal{D}^{U_{\Sigma}} \to \mathcal{C}^{U_{\Gamma}}$, with explicit mate transformation $\beta_N = U(U_{\Sigma}(\varepsilon_N)) \circ \eta_{U_{\Gamma}(UN)}$.
\end{theorem}
\begin{proof}
  \leanok
  Obtained by applying Proposition~\ref{prop:lifted_adjunction} to Theorem~\ref{thm:unstable_alpha}.
\end{proof}

\begin{theorem}[Lifted Counit Isomorphism $\alpha_*\alpha^* N \cong N$]
  \label{thm:unstable_counit_iso}
  \lean{ThesisLean.Section3_2.unstablePushforwardCotensorIso, ThesisLean.Section3_2.isIso_unstableLiftedCounit}
  \uses{thm:unstable_lifted_adj, prop:canonical_equalizer}
  \leanok
  Assume $\mathcal{C}^{U_{\Gamma}}$ has equalizers, $F$ preserves the defining equalizer of $\alpha^* N$, and $\varepsilon : FU \xrightarrow{\sim} 1_{\mathcal{D}}$ is an isomorphism. Then the lifted counit $\tilde{\varepsilon}_N : \alpha_*(\alpha^* N) \to N$ is an isomorphism of $U_{\Sigma}$-coalgebras.
\end{theorem}
\begin{proof}
  \leanok
  On cofree objects, $\alpha_*(\cofree_{U_{\Gamma}}(UY)) = \cofree_{U_{\Sigma}}(Y)$ holds definitionally. Since $F$ preserves the equalizer and $\varepsilon$ is an isomorphism, $\alpha_*(\alpha^* N)$ is the equalizer of $(\cofree_{U_{\Sigma}}(\psi_N), \Delta^{\Sigma}_N)$, which is isomorphic to $N$ by Proposition~\ref{prop:canonical_equalizer} via the lifted counit $\tilde{\varepsilon}_N$.
\end{proof}

\begin{proposition}[Excess Bounds and Colimit Decomposition]
  \label{prop:excess_and_colimit}
  \lean{ThesisLean.Section3_2.excess_vn_shift_iff, ThesisLean.Section3_2.excess_ul_action_iff, ThesisLean.Section3_2.prop3_2_2_colimit_iso}
  \uses{thm:unstable_sigma_comonad}
  \leanok
  The unstable excess condition $2p^n < k + (2p^n - 2)$ holds iff $k \ge 3$, and the $b(n)_*$-action excess condition $2p^n < k$ holds iff $k \ge 2p^n + 1$. In this range, left and right $v_n^{-1}$-actions commute with conjugation, and colimit preservation yields $U_{\Gamma}(K(n)_* \otimes M) \cong B(n)_* \otimes_{K(n)_*} U_{\Sigma}(K(n)_* \otimes M)$.
\end{proposition}
\begin{proof}
  \leanok
  Arithmetic bounds verified by linear arithmetic; inversion and colimit isomorphisms verified categorically.
\end{proof}

\begin{proposition}[Unstable Cotensor Identification and Height-$n$ Unit Isomorphism]
  \label{prop:height_n_unit_iso}
  \lean{ThesisLean.Section3_2.prop3_2_3_unitPullbackIso, ThesisLean.Section3_2.prop3_2_4_cotensor_cofree, ThesisLean.Section3_2.liftedUnitIsoOfHeightNConnected}
  \uses{thm:unstable_lifted_adj, prop:canonical_equalizer, prop:excess_and_colimit}
  \leanok
  When $B(n)_* \otimes_{K(n)_*} -$ preserves the canonical equalizer of $U_{\Sigma}$-coalgebras:
  \begin{enumerate}
    \item $\alpha^*(U_{\Sigma}(Y)) \cong U_{\Gamma}(U(Y))$ on cofree coalgebras.
    \item For $k$-connected $M$ with $k \ge 2p^n + 1$, the underlying object of $\alpha^*\alpha_* M$ is isomorphic to $B(n)_* \otimes_{BP_*} M$.
    \item If $M$ is also of height $n$, the lifted unit $\tilde{\sigma}_M : M \to \alpha^*\alpha_* M$ is an isomorphism of $U_{\Gamma}$-coalgebras.
  \end{enumerate}
\end{proposition}
\begin{proof}
  \leanok
  Proved by comparing the defining equalizer of $\alpha^*$ with the image of the canonical equalizer of Proposition~\ref{prop:canonical_equalizer} under $B(n)_* \otimes_{K(n)_*} -$.
\end{proof}

\begin{theorem}[Unstable Change of Rings]
  \label{thm:unstable_change_of_rings}
  \lean{ThesisLean.Section3_2.unstableChangeOfRingsCobarIso, ThesisLean.Section3_2.unstablePullbackCobarLeftContractibleTower}
  \uses{prop:generalized_shapiro, prop:height_n_unit_iso}
  \leanok
  If $M$ is a $U_{\Gamma}$-comodule of height $n$ and connectivity $k \ge 2p^n + 1$, then for any test coalgebra $P \in \mathcal{C}^{U_{\Gamma}}$, there is an isomorphism of coaugmented cosimplicial cobar hom-towers
  \[
    \Hom_{\mathcal{C}^{U_{\Gamma}}}(P, C^{\bullet}_{U_{\Gamma}}(M)) \;\cong\; \Hom_{\mathcal{D}^{U_{\Sigma}}}(\alpha_* P, C^{\bullet}_{U_{\Sigma}}(\alpha_* M)),
  \]
  and $\alpha^*(C^{\bullet}_{U_{\Sigma}}(\alpha_* M))$ has a left-contractible hom-tower into every cofree $U_{\Gamma}$-coalgebra. Consequently, $\Ext_{U_{\Gamma}}(BP_*, M) \cong \Ext_{U_{\Sigma}}(K(n)_*, K(n)_* \otimes_{BP_*} M)$.
\end{theorem}
\begin{proof}
  \leanok
  Combines the cobar isomorphism induced by the lifted unit isomorphism $\tilde{\sigma}_M : M \xrightarrow{\sim} \alpha^*\alpha_* M$ of Proposition~\ref{prop:height_n_unit_iso} with the Generalized Shapiro's Lemma (Proposition~\ref{prop:generalized_shapiro}).
\end{proof}
